\section{Комбинаторы парсеров для поиска путей с КС ограничениям}
	\tikzsetfigurename{CFPQ_Combinators_}

Основные сведения о комбинаторах парсеров вынесены в раздел~\ref{sec:Combinators}.
В данном разделе мы рассматриваем применение этой техники для задачи поиска путей с контекстно-свободными ограничениями в графах (определение~\ref{def:FLPQ}).

Использование комбинаторов для CFPQ решает сразу две задачи.
Во-первых, оно обеспечивает прозрачную интеграцию языка запросов в язык общего назначения~\sidecite{Verbitskaia:2018:PCC:3241653.3241655}: запросы являются обычными функциями, и компилятор языка проверяет их корректность статически.
Это устраняет проблемы подходов на основе строкового встраивания предметно-ориентированных языков (string-embedded DSL) и ORM/LINQ, в которых декомпозиция и переиспользование подзапросов затруднены.
Во-вторых, комбинаторы предоставляют естественный механизм композиции запросов и позволяют строить параметризованные шаблоны типовых запросов (например, запросов вида <<same generation>>~\sidecite{CFGonRDF}).

Идея использовать комбинаторы для навигации по графам не нова; первая реализация была предложена в библиотеке Trails~\sidecite{Kroni:2013:PGA:2489837.2489844}.
Однако Trails, как и большинство библиотек комбинаторов парсеров, не поддерживает леворекурсивные грамматики и испытывает трудности при обработке циклов во входном графе.
Библиотека Meerkat~\sidecite{izmaylova2016practical} решает эти проблемы за счёт идей обобщённого LL-алгоритма (GLL)~\sidecite{Scott:2010:GP:1860132.1860320}: мемоизация~\sidecite{Johnson1995memoization} и передача продолжений (continuation-passing style) позволяют обрабатывать произвольные КС-грамматики, включая леворекурсивные и неоднозначные.
На основе Meerkat была построена библиотека комбинаторов для графовых запросов~\sidecite{Verbitskaia:2018:PCC:3241653.3241655}, которая и рассматривается далее.

\subsection{Описание алгоритма}
\label{subsec:CFPQ_Combinators_description}

В основе подхода лежит следующее наблюдение.
Классический синтаксический анализатор обрабатывает линейный вход~--- строку, которую можно рассматривать как линейный ориентированный граф: символы метят рёбра между соседними позициями.
Анализатор перемещает указатель с позиции $i$ на позицию $i+1$ и создаёт новое состояние, если токен между позициями $i$ и $i+1$ соответствует правилу грамматики.
При переходе к графу общего вида из текущей вершины может исходить несколько рёбер с подходящей меткой, что приводит к порождению нескольких новых состояний.
Обобщённые алгоритмы синтаксического анализа (GLL и GLR) спроектированы для эффективной обработки множественных состояний и потому естественно обобщаются на графы~\sidecite{Grigorev:2017:CPQ:3166094.3166104}.

\subsubsection{Обобщённый интерфейс входа}

Ключевая идея предлагаемого подхода~--- абстрагирование от конкретного представления графа.
Для работы комбинаторов достаточно двух операций, образующих интерфейс \texttt{Input} с параметрами типов меток рёбер $L$ и меток вершин $N$:
\begin{itemize}
    \item \texttt{filterEdges} — по заданной вершине и предикату на метках рёбер возвращает множество исходящих рёбер (метка и целевая вершина);
    \item \texttt{checkNode}~--- проверяет, удовлетворяет ли заданная вершина предикату, и возвращает её метку.
\end{itemize}
Благодаря такому интерфейсу библиотека комбинаторов не зависит от источника данных.
Поддерживаются графы, загруженные в оперативную память, графовая база данных Neo4j и даже линейный вход (обычные строки).

\subsubsection{Базовые комбинаторы}

Библиотека предоставляет набор примитивных комбинаторов и операций для построения более сложных запросов.
Два базовых строительных блока~--- комбинаторы для работы с вершинами и рёбрами:
\begin{itemize}
    \item $V(predicate)$~--- проверяет, что текущая вершина удовлетворяет предикату $predicate$;
    \item $E(predicate)$~--- проверяет, что из текущей вершины исходит ребро, метка которого удовлетворяет предикату $predicate$.
\end{itemize}

На основе этих примитивов строятся более сложные запросы с помощью операций, перечисленных в таблице~\ref{table:CFPQ_Combinators}.
Макрос \texttt{syn} (от \emph{syntax}) преобразует строковый литерал или композицию комбинаторов в рекурсивный нетерминал; например, запись \texttt{val S = syn("a" $\sim$ S.? $\sim$ "b")} задаёт язык $\{a^n b^n \mid n \geq 1\}$.

\begin{table}[h]
    \centering
    \caption{Основные комбинаторы}
    \label{table:CFPQ_Combinators}
    \begin{tabular}{l|l}
        \hline
        Комбинатор & Описание \\
        \hline
        $a \sim b$  & последовательное применение $a$, затем $b$ \\
        $a \mid b$  & альтернатива: $a$ или $b$ \\
        $a?$        & опциональный парсер: $a$ или ничего \\
        $a*$        & повторение нуль или более раз \\
        $a+$        & повторение один или более раз \\
        $a \mathbin{\texttt{\textasciicircum}} f$ & применить функцию $f$ к результату токена $a$ \\
        $a \mathbin{\texttt{\&}} f$ & применить функцию $f$ к результату композиции $a$ \\
        \hline
    \end{tabular}
\end{table}

\subsubsection{Семантические действия}

Для извлечения данных из результатов запроса библиотека предоставляет механизм семантических действий~--- аналог семантических действий в классических генераторах парсеров.
Различают две разновидности:

\begin{itemize}
    \item \texttt{\textasciicircum} применяется к терминальным комбинаторам ($V$ и $E$) и позволяет извлечь свойства токена (например, имя вершины или длину пути);
    \item \texttt{\&} применяется к композиции комбинаторов и позволяет собрать и обработать результаты подзапросов (например, сформировать пару из начальной и конечной вершины пути).
\end{itemize}

Семантические действия вычисляются снизу вверх по дереву разбора (SPPF): сначала обрабатываются действия для дочерних узлов, затем их результаты передаются в действие текущего узла.
Поскольку SPPF может содержать неоднозначные узлы (соответствующие нескольким различным путям), библиотека предоставляет механизм ленивого обхода SPPF в ширину для извлечения деревьев разбора.

Функция \texttt{executeQuery} запускает разбор входного графа из всех вершин, строит SPPF, извлекает все деревья разбора, применяет семантические действия и возвращает ленивый поток результатов.

\subsubsection{SPPF как представление результата}

В основе библиотеки Meerkat лежит обобщённый LL-алгоритм, который строит сжатое представление леса разбора~--- SPPF (раздел~\ref{sec:GLR}).
Как было показано ранее~\sidecite{Grigorev:2017:CPQ:3166094.3166104}, SPPF является подходящим конечным представлением потенциально бесконечного множества путей.
Все пути могут быть извлечены из SPPF в виде соответствующих деревьев вывода.

\subsubsection{Обработка циклов и левой рекурсии}

Мемоизация, используемая в Meerkat, решает две ключевые проблемы, возникающие при обработке графов.
Во-первых, разбор из заданного состояния в заданной вершине выполняется только один раз; поскольку число состояний и вершин конечно, разбор гарантированно завершается, и бесконечные циклы при обходе циклов графа невозможны.
Во-вторых, появление нового состояния разбора в вершине, которая уже была обработана, запускает пересчёт всех зависимых от него путей.
Благодаря мемоизации каждый вывод каждого подпути анализируется ровно один раз, что не приводит к существенным накладным расходам.

Этот же механизм мемоизации и передачи продолжений (continuation-passing style) позволяет Meerkat обрабатывать леворекурсивные комбинаторы~\sidecite{izmaylova2016practical}, что является общей проблемой подхода комбинаторов парсеров.

\subsection{Примеры}
\label{subsec:CFPQ_Combinators_examples}

\subsubsection{Запрос same generation}

Рассмотрим классический контекстно-свободный запрос <<same generation>>~\sidecite{CFGonRDF}, полезный для анализа иерархий в RDF-графах.
Пусть граф содержит отношения двух типов: \texttt{subClassOf} и \texttt{type}, а также обратные к ним \texttt{subClassOf}$^{-1}$ и \texttt{type}$^{-1}$.
Задача~--- найти вершины, находящиеся на одном уровне иерархии, то есть достижимые путями, в которых каждому шагу <<вверх>> соответствует шаг <<вниз>> по парному отношению.

Грамматика для этого запроса на двух парах отношений имеет вид:
\[
    \mathit{qSameGen} \to \mathit{R}_0^{-1}\ \mathit{qSameGen?}\ \mathit{R}_0 \mid \mathit{R}_1^{-1}\ \mathit{qSameGen?}\ \mathit{R}_1
\]
где $\mathit{R}_0 = \text{\texttt{subClassOf}}$, $\mathit{R}_1 = \text{\texttt{type}}$, а $\mathit{R}_i^{-1}$ обозначает обратное отношение.
Средствами комбинаторов этот запрос записывается следующим образом:

\begin{verbatim}
val qSameGen = syn(
    inE(_.label == "subClassOf") ~ qSameGen.? ~ outE(_.label == "subClassOf") |
    inE(_.label == "type") ~ qSameGen.? ~ outE(_.label == "type"))
\end{verbatim}

% TODO: Здесь было бы хорошо привести рисунок графа, на котором демонстрируется запрос.
% \mytodo{Создать рисунок графа для примера same generation.}

\subsubsection{Параметризованные комбинаторы и композиция}

Грамматики для запросов <<same generation>> над разными наборами пар отношений структурно идентичны и различаются лишь конкретными отношениями.
Средствами комбинаторов это обобщение выражается естественно: функция \texttt{sameGen} принимает список пар <<скобок>> и строит соответствующий запрос.

Вспомогательная функция \texttt{reduceChoice} сводит список подзапросов к одному с помощью операции альтернативы:

\begin{verbatim}
def reduceChoice(qs: List[Nonterminal]) = qs match {
    case x :: Nil => x
    case x :: y :: qs => syn(qs.foldLeft(x | y)(_ | _))
}
\end{verbatim}

Тогда параметризованный комбинатор \texttt{sameGen} определяется как:

\begin{verbatim}
def sameGen(brs: List[(Nonterminal, Nonterminal)]) =
    reduceChoice(brs.map { case (lbr, rbr) =>
        syn(lbr ~ sameGen(brs).? ~ rbr) })
\end{verbatim}

Теперь запросы над разными наборами отношений конструируются простым применением \texttt{sameGen} к нужному списку пар, без дублирования кода.

\subsubsection{Семантические действия и типобезопасность}

Продолжим пример запроса <<same generation>>.
Пусть требуется не только найти достижимые вершины, но и вычислить длину каждого пути.
Добавим семантические действия для аккумулирования длины:

\begin{verbatim}
def sameGen(brs: List[(_, _)]) =
    reduceChoice(brs.map { case (lbr, rbr) =>
        syn((lbr ~ (sameGen(brs).?) ~ rbr) & {
            case _ ~ Nil ~ _       => 2
            case _ ~ (len: Int) ~ _ => len + 2
        })
    })
\end{verbatim}

Здесь \texttt{\&} применяется к результату последовательности из трёх элементов.
Если средний элемент отсутствует (путь нулевой длины между <<скобками>>), длина пары скобок равна 2; иначе к накопленной длине добавляются два ребра текущей пары.

Поскольку запросы являются обычными функциями языка Scala, компилятор статически проверяет корректность типов.
Например, если в семантическом действии вместо целого числа попытаться использовать строку, компилятор укажет на ошибку типа.
Кроме того, IDE предоставляет подсветку синтаксиса, навигацию по коду и автодополнение без дополнительных настроек.

\begin{example}
    Рассмотрим запрос к онтологии \emph{Geospecies}, содержащей биологическую иерархию.
    Для стартовой вершины $v$ запрос на основе комбинатора \texttt{sameGen} с отношением \texttt{skos\_\_narrowerTransitive} возвращает множество объектов одного уровня иерархии.
    Одновременно вычисляется длина соответствующих путей.
    Эксперименты на реальных RDF-графах показывают, что время выполнения и потребление памяти линейны по размеру результирующего множества путей.
\end{example}

\subsection{Свойства алгоритма}
\label{subsec:CFPQ_Combinators_properties}

Опишем основные свойства подхода.

\begin{theorem}[Корректность]
    Множество путей, построенное библиотекой комбинаторов для запроса, заданного КС-грамматикой $G$ и графа $\mbfscrG$, совпадает с множеством путей задачи CFPQ для $G$ и $\mbfscrG$.
\end{theorem}

\begin{proof}
    Корректность следует из корректности GLL-алгоритма~\sidecite{Scott:2010:GP:1860132.1860320} и его обобщения на графы~\sidecite{Grigorev:2017:CPQ:3166094.3166104}.
    Библиотека Meerkat реализует обобщённый LL-алгоритм, а обобщение на графы достигается заменой функции перехода <<следующая позиция>> на функцию \texttt{filterEdges}, возвращающую множество следующих вершин.
    Корректность этого обобщения показана в~\sidecite{Verbitskaia:2018:PCC:3241653.3241655}.
\end{proof}

\noindent\textbf{Выразительная мощность.}
Подход поддерживает произвольные КС-грамматики, включая леворекурсивные и неоднозначные.
Это обеспечивается мемоизацией и передачей продолжений~\sidecite{izmaylova2016practical,Johnson1995memoization}.

\medskip
\noindent\textbf{Обработка циклов.}
Мемоизация гарантирует завершение разбора на графах с циклами: каждое состояние разбора в каждой вершине обрабатывается не более одного раза, а число состояний конечно.

\medskip
\noindent\textbf{Сложность.}
В худшем случае временная и пространственная сложность кубические относительно размера графа, что унаследовано от GLL-алгоритма~\sidecite{Scott:2010:GP:1860132.1860320}.
На практике для задачи с одним источником (single-source CFPQ) время выполнения часто линейно по размеру результата, что подтверждается экспериментальными данными~\sidecite{Verbitskaia:2018:PCC:3241653.3241655}.

\medskip
\noindent\textbf{Представление результата.}
SPPF (раздел~\ref{sec:GLR}) является конечной структурой данных, представляющей потенциально бесконечное множество путей.
Из SPPF могут быть лениво извлечены все деревья вывода с применением семантических действий.

\medskip
\noindent\textbf{Сравнение с другими подходами.}
В отличие от матричных алгоритмов (раздел~\ref{sec:MatrixCFPQ}), вычисляющих достижимость между всеми парами вершин, комбинаторы естественно поддерживают семантику с одним источником и не требуют приведения грамматики к нормальной форме Хомского, как подходы на основе CYK (раздел~\ref{sec:CYK}), GLL (раздел~\ref{sec:CFPQ_GLL}) и GLR (раздел~\ref{sec:GLR}).
Дополнительным преимуществом является возможность задания пользовательских семантических действий и параметризованных шаблонов запросов, обеспечиваемая интеграцией в язык общего назначения.
